\documentclass{article}
\usepackage{amsmath, amssymb, amsthm}
\usepackage{hyperref}
\usepackage{geometry}
\geometry{margin=1in}

\title{Erd\H{o}s Problem \#332}
\date{Last updated: 2025-08-31}
\author{Source: erdosproblems.com}

\begin{document}
\maketitle

\noindent\textbf{Status:} OPEN \quad \textbf{No prize}

\noindent\textbf{Tags:} number theory

\medskip
\noindent\textbf{Problem Statement:}

Let $A\subseteq \mathbb{N}$ and $D(A)$ be the set of those numbers which occur infinitely often as $a_1-a_2$ with $a_1,a_2\in A$. What conditions on $A$ are sufficient to ensure $D(A)$ has bounded gaps?

\bigskip
\noindent\textbf{Remarks:}

Prikry, Tijdeman, Stewart, and others (see the survey articles \cite{St78} and \cite{Ti79}) have shown that a sufficient condition is that $A$ has positive density.

One can also ask what conditions are sufficient for $D(A)$ to have positive density, or for $\sum_{d\in D(A)}\frac{1}{d}=\infty$, or even just $D(A)\neq\emptyset$.

\bigskip
\noindent\textbf{References:}

[St78] Stewart, Cam L., On difference sets of sets of integers. S\'{e}minaire Delange-Pisot-Poitou, 19e ann\'{e}e:
1977/78, Th\'{e}orie des nombres, Fasc. 1 (1978), Exp. No. 5, 8.
 
 [Ti79] Tijdeman, R., Distance sets of sequences of integers. Proceedings, Bicentennial Congress Wiskundig
Genootschap (Vrije Univ., Amsterdam, 1978), Part
II (1979), 405-415.

\bigskip
\noindent\small{Source: \url{https://www.erdosproblems.com/332}}
\end{document}
